% =============================================================================
% Configuracion general del documento.
% Si tu plantilla oficial de la ETSIIT viene con su propio preamble, sustituye
% este archivo por el suyo y deja tal cual el resto (capitulos y main.tex).
% =============================================================================


\usepackage{float}
\usepackage{subcaption}
\usepackage{placeins}

% ----- Codificacion e idioma
\usepackage[utf8]{inputenc}
\usepackage[T1]{fontenc}
\usepackage[main=spanish,english]{babel}
\usepackage{csquotes}

% ----- Margenes (la ETSIIT recomienda ~2,5 cm; ajusta si tu tutor lo pide)
\usepackage[a4paper, top=2.5cm, bottom=2.5cm, inner=3cm, outer=2.5cm]{geometry}

% ----- Tipografias
\usepackage{lmodern}
\usepackage{microtype}

% ----- Matematicas y simbolos
\usepackage{amsmath,amssymb,amsthm}
\usepackage{mathtools}
\usepackage{siunitx}
\sisetup{
  output-decimal-marker={,},
  detect-all,
  per-mode=symbol,
  group-separator={\,}
}

% ----- Graficos y figuras
\usepackage{graphicx}
\graphicspath{{figuras/}}
\usepackage{float}
\usepackage{caption}
\usepackage{subcaption}
\captionsetup[figure]{font=small,labelfont=bf,name=Figura}
\captionsetup[table]{font=small,labelfont=bf,name=Tabla}

% ----- Colores
\usepackage{xcolor}
\definecolor{ugrazul}{HTML}{1F3864}
\definecolor{ugrgris}{HTML}{595959}

% ----- Tablas
\usepackage{booktabs}
\usepackage{multirow}
\usepackage{array}
\usepackage{tabularx}

% ----- Codigo fuente (listings)
\usepackage{listings}
\lstset{
  basicstyle=\ttfamily\footnotesize,
  keywordstyle=\color{ugrazul}\bfseries,
  commentstyle=\color{ugrgris}\itshape,
  stringstyle=\color{teal},
  numbers=left,
  numberstyle=\tiny\color{ugrgris},
  stepnumber=1,
  numbersep=8pt,
  showspaces=false,
  showstringspaces=false,
  showtabs=false,
  frame=single,
  framerule=0.4pt,
  rulecolor=\color{ugrgris},
  tabsize=4,
  captionpos=b,
  breaklines=true,
  breakatwhitespace=true,
  inputencoding=utf8,
  extendedchars=true,
  literate={\'a}{{\'a}}1 {\'e}{{\'e}}1 {\'i}{{\'i}}1 {\'o}{{\'o}}1 {\'u}{{\'u}}1
           {\~n}{{\~n}}1 {\~N}{{\~N}}1 {\'A}{{\'A}}1 {\'E}{{\'E}}1 {\'I}{{\'I}}1
           {\'O}{{\'O}}1 {\'U}{{\'U}}1
}
% Estilo predefinido para Python
\lstdefinestyle{python}{language=Python}
\lstdefinestyle{bash}{language=bash}

% ----- Algoritmos (opcional, util para describir el pipeline)
\usepackage{algorithm}
\usepackage{algpseudocode}

% ----- Encabezados y pies de pagina
\usepackage{fancyhdr}
\setlength{\headheight}{14.5pt}
\pagestyle{fancy}
\fancyhf{}
\fancyhead[LE]{\nouppercase\leftmark}
\fancyhead[RO]{\nouppercase\rightmark}
\fancyfoot[LE,RO]{\thepage}
\renewcommand{\headrulewidth}{0.4pt}
\renewcommand{\footrulewidth}{0pt}

% ----- Bibliografia con biblatex + biber (estandar moderno)
\usepackage[backend=biber,style=ieee,sorting=none,maxbibnames=10]{biblatex}
\addbibresource{bibliografia.bib}

% ----- Hipervinculos y referencias cruzadas
\usepackage[hidelinks]{hyperref}
\usepackage[spanish]{cleveref}
\crefname{chapter}{capítulo}{capítulos}
\Crefname{chapter}{Capítulo}{Capítulos}
\crefname{section}{sección}{secciones}
\Crefname{section}{Sección}{Secciones}
\crefname{figure}{figura}{figuras}
\Crefname{figure}{Figura}{Figuras}
\crefname{table}{tabla}{tablas}
\Crefname{table}{Tabla}{Tablas}
\crefname{equation}{ecuación}{ecuaciones}
\Crefname{equation}{Ecuación}{Ecuaciones}

% ----- Glosario / acronimos (paquete simple)
\usepackage[printonlyused,withpage]{acronym}

% ----- Notas: usar \todo{texto} durante la redaccion
\usepackage[textsize=small,color=yellow!40]{todonotes}

% ----- Reduce huerfanas/viudas y mejora el flujo del texto
\widowpenalty=10000
\clubpenalty=10000
